% Конспект-методичка лекции 11. Сборка из корня:
% python3 11/images/overview_numerics.py
% python3 11/images/overview_pointing.py
% python3 11/images/overview_timing.py
% python3 11/images/overview_diagrams.py
% python3 11/images/galileo_case.py
% latexmk 11/practical-operations.tex
\documentclass[a4paper,11pt]{article}
\usepackage{iftex}
\RequirePDFTeX
\usepackage{cmap}
\usepackage[T2A]{fontenc}
\usepackage[utf8]{inputenc}
\usepackage[english,russian]{babel}
\usepackage{PTSerif}
\usepackage{amsmath,amssymb,amsthm,mathtools}
\usepackage[a4paper,left=20mm,right=20mm,top=22mm,bottom=23mm,headheight=14pt]{geometry}
\usepackage{microtype}
\usepackage{graphicx,xcolor,tikz,booktabs,tabularx,array}
\usetikzlibrary{arrows.meta,calc,positioning}
\usepackage{enumitem,fancyhdr,needspace,float,placeins}
\usepackage[most]{tcolorbox}
\definecolor{accent}{HTML}{187AA5}
\definecolor{muted}{HTML}{5A6166}
\definecolor{figink}{HTML}{25323F}
\definecolor{figtan}{HTML}{087F78}
\definecolor{figorange}{HTML}{B57714}
\definecolor{figpurple}{HTML}{79539C}
\newtcolorbox{explanation}[1]{enhanced,
 colback=accent!3!white,colframe=accent!80!black,
 boxrule=.5pt,toprule=1.2pt,arc=1mm,
 left=3mm,right=3mm,top=3mm,bottom=2mm,
 before skip=9pt,after skip=9pt,
 title={#1},fonttitle=\bfseries,coltitle=accent!70!black,
 attach boxed title to top left={xshift=3mm,yshift=-2mm},
 boxed title style={colback=white,frame hidden,boxrule=0pt,
 left=1.5mm,right=1.5mm,top=.4mm,bottom=.4mm}}
\newtcolorbox{remarkbox}{enhanced,
 colback=orange!4!white,colframe=orange!65!black,
 boxrule=.5pt,toprule=1.2pt,arc=1mm,
 left=3mm,right=3mm,top=2mm,bottom=2mm,
 before skip=9pt,after skip=9pt}
\usepackage{hyperref}
\usepackage{bookmark}
\hypersetup{unicode=true,colorlinks=true,linkcolor=accent,citecolor=accent,urlcolor=accent,
 pdftitle={Практические аспекты орбитальных операций},
 pdfauthor={Материалы курса «Введение в маневрирование КА»},
 pdfsubject={Обзор ошибок, двигательной установки, нештатных ситуаций и цикла орбитальных операций}}
\newcommand{\dd}{\mathop{}\!\mathrm d}
\newcommand{\vect}[1]{\boldsymbol{#1}}
\newcommand{\norm}[1]{\left\lVert#1\right\rVert}
\numberwithin{equation}{section}
\setcounter{tocdepth}{2}
\setcounter{secnumdepth}{2}
\setlength{\parindent}{1.3em}
\setlength{\parskip}{2pt}
\linespread{1.055}
\setlist{itemsep=2pt,topsep=4pt,parsep=0pt}
\setlength{\emergencystretch}{1.5em}
\clubpenalty=10000
\widowpenalty=10000
\brokenpenalty=10000
\displaywidowpenalty=10000
\allowdisplaybreaks[1]
\setcounter{topnumber}{2}
\renewcommand{\topfraction}{.88}
\renewcommand{\textfraction}{.10}
\renewcommand{\floatpagefraction}{.82}
\makeatletter
\setlength{\@fptop}{0pt}
\makeatother
\pagestyle{fancy}
\fancyhf{}
\fancyhead[L]{\small\color{muted}Практические аспекты орбитальных операций}
\fancyfoot[C]{\small\thepage}
\renewcommand{\headrulewidth}{.3pt}
\newcommand{\coursefigure}[3][.98]{%
\begin{figure}[!htbp]\centering
\includegraphics[width=#1\textwidth]{images/#2.pdf}
\caption{#3}\label{fig:#2}\end{figure}}
% Generated by overview_numerics.py; all values belong to stated teaching models.
\newcommand{\OvInjectionNominalDV}{1453{,}076}
\newcommand{\OvInjectionMinDV}{1416{,}532}
\newcommand{\OvInjectionMaxDV}{1492{,}240}
\newcommand{\OvFuelNominal}{66{,}839}
\newcommand{\OvFuelLowIsp}{74{,}126}
\newcommand{\OvDurationNominal}{436{,}976}
\newcommand{\OvDurationLowThrust}{655{,}464}
\newcommand{\OvDurationNominalMinutes}{7{,}283}
\newcommand{\OvDurationLowThrustMinutes}{10{,}924}
\newcommand{\OvMinimumThrust}{327{,}732}
\newcommand{\OvLimitedDV}{91{,}405}
\newcommand{\OvMixtureResidual}{5{,}556}

% Generated by overview_timing.py; do not edit numeric values manually.
\newcommand{\TimingInitialPeriod}{9215{,}022}
\newcommand{\TimingInitialHours}{2{,}560}
\newcommand{\TimingTargetRp}{7463{,}631}
\newcommand{\TimingDelayShortCost}{0{,}186}
\newcommand{\TimingDelayMediumCost}{4{,}651}
\newcommand{\TimingDelayLongCost}{10{,}461}
\newcommand{\TimingDelayShortPeriDV}{-0{,}087793}
\newcommand{\TimingDelayShortApoDV}{0{,}098305}
\newcommand{\TimingDelayShortFinalTime}{9538{,}779}
\newcommand{\TimingBurnDuration}{436{,}976}
\newcommand{\TimingAbortHalfDV}{49{,}575}
\newcommand{\TimingAbortHalfRp}{7225{,}658}
\newcommand{\TimingAbortFullRp}{7461{,}583}
\newcommand{\TimingAbortFullRa}{12002{,}950}
\newcommand{\TimingAbortFullOmega}{-0{,}788}
\newcommand{\TimingDelayLongOmega}{-3{,}120}

\begin{document}
\begin{titlepage}
\centering
{\large Московский физико-технический институт\par}
\vspace{8mm}
{\large Курс «Введение в маневрирование КА»\par}
\vspace{30mm}
{\color{accent}\rule{28mm}{.8pt}\par}
\vspace{9mm}
{\LARGE\bfseries Практические аспекты\\[3mm]орбитальных операций\par}
\vspace{11mm}
{\large Методические материалы к лекции 11\par}
\vspace{8mm}
{Анализ ошибок и возможности их исправления.\\
Двигательная установка и нештатные ситуации.\\
Планирование, подготовка и контроль манёвров.\par}
\vfill
{Рабочая редакция\par}
\vspace{3mm}
{17 ноября 2026 года\par}
\end{titlepage}
\setcounter{page}{2}
\begingroup\small
\setlength{\parskip}{0pt}
\tableofcontents
\endgroup
\clearpage

\section*{Предмет рассмотрения и обозначения}
\addcontentsline{toc}{section}{Предмет рассмотрения и обозначения}

Практическая орбитальная операция связывает баллистический расчёт
с возможностями аппарата, измерениями и действиями команды управления.
Даже правильно найденный манёвр может потребовать изменения:
начальная орбита оказалась другой, тяга ниже ожидаемой,
подходящее окно пропущено или двигатель прекратил работу раньше времени.

Рассмотрим основные классы таких ситуаций, их последствия для
рабочего тела и времени полёта, а затем полный цикл подготовки
и исполнения манёвра. Основой служит глава~15
DiPrinzio~\cite{diprinzio}. Учебные расчёты дополняют два реальных
примера: исправление орбиты Galileo и восстановление миссии Akatsuki.

В схемах используются плановое состояние, фактическое исполнение
и варианты продолжения. Численные параметры учебных аппаратов
задаются отдельно; они не являются характеристиками Galileo или Akatsuki.

\begin{center}
\renewcommand{\arraystretch}{1.12}
\begin{tabularx}{\textwidth}{@{}>{$}l<{$}>{\raggedright\arraybackslash}X@{}}
\toprule
\multicolumn{1}{@{}l}{Обозначение}&Смысл\\\midrule
\vect R,\vect V & Инерциальные положение и скорость КА.\\
r_p,r_a & Радиусы перигея и апогея, измеряемые от центра Земли.\\
\Delta\vect V & Вектор мгновенного приращения скорости.\\
\Delta V,J & Характеристическая скорость одного манёвра и сумма
характеристических скоростей последовательности.
Для мгновенного манёвра $\Delta V=\norm{\Delta\vect V}$.\\
F,m,I_{\mathrm{sp}} & Модуль тяги, текущая масса и удельный импульс двигателя.\\
\bottomrule
\end{tabularx}
\end{center}
Остальные обозначения вводятся по месту. Углы в формулах измеряются
в радианах; градусы в численных примерах отмечены явно.

\section{Анализ ошибок и планирование миссии}
\label{sec:planning}

Номинальная программа отвечает на вопрос, как выполнить переход
при принятых исходных данных. Анализ ошибок проверяет, какие изменения
этой программы потребуются при отклонениях и останется ли продолжение
допустимым.

\subsection{Типы ошибок и их последствия}
\label{sec:classification}

Для орбитального перехода выделим пять основных типов сценариев:
\begin{enumerate}
\item \textbf{Ошибки выведения.}
Орбита после отделения от ракеты-носителя отличается от номинальной;
меняются начальные условия последующего перехода.
\item \textbf{Ошибки наведения.}
Фактическое направление управляющего воздействия отличается
от требуемого.
\item \textbf{Ошибки работы двигательной установки.}
Тяга, удельный импульс, расход компонентов или доступная конфигурация
двигателей отличаются от принятых в расчёте.
\item \textbf{Опоздание манёвра.}
Воздействие начинается позже запланированного момента,
при изменившемся положении и скорости КА.
\item \textbf{Пропуск манёвра.}
Воздействие не выполнено в предусмотренном окне;
необходимо выбрать другое продолжение.
\end{enumerate}
Отдельно рассматривается \textbf{аварийное прекращение манёвра}:
часть воздействия уже получена, после чего работа двигателя
прервалась. Эти ситуации могут сочетаться.

У отклонения есть несколько разных последствий. Исправление
может потребовать дополнительного рабочего тела или увеличить время
перехода. Может измениться доступность сеансов связи, освещённость,
режим полезной нагрузки или возможность работы датчиков.
Наконец, прежняя цель может стать недостижимой с оставшимися
двигателями и ресурсами.

Большая ошибка не всегда критична, если предусмотрен достаточный
резерв управления. Небольшое отклонение может оказаться критичным,
если после него исчезает последнее допустимое окно.
Поэтому оценивают \emph{возможность продолжения миссии},
а не только величину ошибки.

\subsection{Что проверяют до запуска и в полёте}
\label{sec:mission-stages}

До запуска рассчитывают номинальный переход и варианты
для ожидаемых отклонений. Например, изменяют исходную орбиту,
уровень тяги или доступный двигатель и заново строят последовательность.
Так получают требования к резерву рабочего тела, времени
и возможностям управления.

После выведения исходные предположения заменяются оценкой
фактического состояния КА. После каждого манёвра уточняются его
результат и модель исполнения. Эти этапы и возможные решения
показаны на рисунке~\ref{fig:overview-planning}.

\coursefigure[.94]{overview-planning}{Анализ ошибок на разных этапах
миссии. От полученных данных переходят к выбору допустимого
продолжения: коррекции, переносу времени, изменению цели
или отказу от невыполнимого варианта. Схема качественная;
численные границы допустимости задаются конкретной задачей.}

При анализе используют отдельные характерные сценарии,
выборки случайных ошибок и расчёты предельных случаев.
Случайная выборка показывает разброс результатов в принятой модели;
она не заменяет проверку отказа двигателя или пропуска окна.
Для каждого нового плана вновь проверяют конечные условия,
ограничения на всём переходе и остающийся ресурс.
\FloatBarrier

\section{Ошибки выведения}
\label{sec:injection}

Ошибка выведения изменяет начальные условия ещё до первого
рассчитанного манёвра КА. Нужно определить, можно ли достичь
целевой орбиты с полученного состояния и сколько это потребует
рабочего тела и времени.

\subsection{Семейство начальных орбит}

Отклонения могут относиться к размерам и форме орбиты,
положению её плоскости и положению КА на орбите.
Они описываются оценками состояния и их неопределённостью.
Для каждой рассматриваемой начальной орбиты рассчитывается
свой вариант дальнейшего маневрирования.

На рисунке~\ref{fig:overview-injection} приведена учебная выборка
для перехода с геопереходной орбиты на геостационарную.
Общий апогей равен $r_a=42164$~км, номинальный перигей ---
$r_p=7000$~км, номинальное наклонение к экватору ---
$i=6^\circ$. Апогей находится на экваториальной линии узлов,
поэтому в нём можно совместить увеличение скорости
и поворот плоскости одним импульсом.

\coursefigure[.94]{overview-injection}{Учебная выборка орбит выведения
и требуемый апогейный импульс. Случайно изменены радиус перигея
и наклонение; общий апогей находится на линии узлов.
Цвет различает величину импульса для каждой начальной орбиты.
Это отдельная модель разброса, а не данные ракеты-носителя.}

Пусть $V_a$ --- апогейная скорость на исходной орбите,
а $V_c=\sqrt{\mu/r_a}$ --- скорость на целевой круговой орбите.
Гравитационный параметр Земли
$\mu=398600.4418$~км$^3$/с$^2$.
В принятой геометрии угол между скоростями равен $i$.
Из треугольника скоростей получаем
\begin{equation}
\Delta V
=\sqrt{V_a^2+V_c^2-2V_aV_c\cos i}.
\label{eq:injection-dv}
\end{equation}
Для номинальных данных это $\OvInjectionNominalDV$~м/с.
В выборке изменённые радиусы перигея и наклонения
приводят к разным затратам на тот же переход.

Выборка содержит 400 состояний с независимыми нормальными
отклонениями: среднеквадратические значения равны $120$~км
для $r_p$ и $0.5^\circ$ для $i$. Эти параметры назначены
для иллюстрации. Если требуемое управление выходит за доступные
возможности, рассматривают другой срок, другую последовательность
или другую целевую орбиту.

\subsection{Реальный пример: Galileo}
\label{sec:galileo}

Пятый и шестой спутники Galileo после запуска в августе 2014~года
оказались на нештатных эллиптических орбитах. Для пятого аппарата
провели 11~манёвров за 17~дней, повысив перигей
и уменьшив эксцентриситет~\cite{esa-recovery}.

Исправленная орбита осталась эллиптической.
В опубликованных данных ESA указаны округлённые высоты
перигея и апогея $17200$ и $25900$~км;
штатная круговая орбита имеет высоту $23222$~км~\cite{esa-facts}.
Их размеры сопоставлены на рисунке~\ref{fig:galileo-recovery}.
Здесь приведены \emph{высоты над поверхностью Земли},
тогда как $r_p,r_a$ в формулах измеряются от её центра.

\coursefigure[.92]{galileo-recovery}{Орбиты Galileo по опубликованным
высотам ESA: после выведения, после исправления и штатная круговая.
Слева их ориентация условно совмещена, справа показаны диапазоны высот.
Это сравнение размеров орбит, а не восстановление траектории
реальных манёвров.}

Повышение перигея уменьшило радиационное воздействие
и обеспечило непрерывную работу датчика Земли.
Так появилась возможность использовать новую орбиту для работы
оборудования и испытаний навигационной нагрузки~\cite{esa-recovery}.
Пример показывает, почему приемлемую цель иногда выбирают
с учётом функций аппарата, а не только исходных орбитальных требований.
\FloatBarrier

\section{Ошибки наведения}
\label{sec:pointing}

Манёвр должен изменить движение в определённом направлении.
Даже при точном модуле импульса ошибка ориентации создаёт
другие компоненты скорости и другой орбитальный результат.

\subsection{Требуемая, фактическая и оценённая ориентация}

\textit{Требуемая ориентация} направляет тягу согласно расчёту.
\textit{Фактическая ориентация} описывает положение аппарата
во время исполнения. \textit{Оценённая ориентация} определяется
по измерениям и сама имеет ошибку.

Если измеренное направление близко к требуемому, это ещё
не означает точного совпадения фактического направления с планом.
На рисунке~\ref{fig:overview-pointing},~а показано
увеличение концов трёх векторов с общим началом.
Углы $1^\circ$ и $0.7^\circ$ назначены для иллюстрации
фактического и оценённого отклонений.

\coursefigure[.94]{overview-pointing}{Наведение и орбитальный результат.
а) Требуемое, фактическое и оценённое направления; начала векторов
находятся вне увеличенной области. б) Угол между начальной
и полученной плоскостями при одинаковом отклонении трансверсального
импульса к нормали в перигее и апогее. Данные учебные.}

\subsection{Почему важно место манёвра}

Пусть плановый импульс направлен по трансверсали, его модуль
$D=100$~м/с, а фактическое направление отклонено к нормали
на $\theta=1^\circ$. Нормальная компонента равна
\[
\Delta V_n=D\sin\theta=1.7452~\text{м/с}.
\]
Она поворачивает плоскость орбиты.
Для одной и той же исходной геопереходной орбиты
из предыдущего раздела отдельный манёвр в перигее поворачивает
плоскость примерно на $0.0100^\circ$,
а в апогее --- на $0.0574^\circ$
(рис.~\ref{fig:overview-pointing},~б).
В апогее полная скорость меньше, поэтому та же добавленная
нормальная компонента сильнее меняет её направление.

Угол между плоскостями и изменение одного наклонения
в общем случае различаются: поворот может изменить также линию узлов.
При анализе наведения оценивают весь орбитальный результат
и возможность его исправить последующими манёврами.
Ошибка в начале последовательности иногда поглощается пересчётом
оставшегося управления; ошибка последнего манёвра
может потребовать отдельной коррекции.
\FloatBarrier

\section{Ошибки и отказы двигательной установки}
\label{sec:propulsion}

Двигательная установка (ДУ) определяет не только сообщённое
воздействие, но и его длительность, расход рабочего тела
и доступные направления. Разные отклонения её характеристик
требуют разных изменений плана.

\subsection{Эффективность и удельный импульс}
\label{sec:isp}

Удельный импульс $I_{\mathrm{sp}}$ связывает тягу с расходом
рабочего тела. Для манёвра с конечной тягой
характеристическая скорость равна
$\Delta V=\int_{t_b}^{t_e}(F/m)\,\dd t$,
где $t_b,t_e$ --- начало и конец работы двигателя.
Это интеграл модуля тягового ускорения; изменение орбитальной
скорости зависит также от направления тяги и остальных сил.
При постоянной скорости истечения
$c=g_0I_{\mathrm{sp}}$ уравнение Циолковского даёт
\begin{equation}
\Delta V=c\ln\frac{m_0}{m_f},\qquad
\Delta m=m_0\left(1-\exp\left[-\frac{\Delta V}{c}\right]\right).
\label{eq:fuel}
\end{equation}
Здесь $g_0=9.80665$~м/с$^2$ --- стандартное ускорение,
$m_0,m_f$ --- массы до и после манёвра,
$\Delta m$ --- израсходованная масса рабочего тела.

Для начальной массы $m_0=2000$~кг
и требуемой характеристической скорости $\Delta V=100$~м/с
при $I_{\mathrm{sp}}=300$~с требуется
$\OvFuelNominal$~кг рабочего тела.
При $270$~с требуется уже $\OvFuelLowIsp$~кг.
Зависимость показана на рисунке~\ref{fig:overview-propulsion},~а.
Пониженная эффективность уменьшает располагаемый ресурс
дальнейших операций, даже если нужный манёвр удаётся выполнить.

\coursefigure[.94]{overview-propulsion}{Два разных отклонения ДУ
для учебного аппарата массой $2000$~кг и требуемых $100$~м/с.
а) Расход рабочего тела при изменении удельного импульса.
б) Длительность при изменении тяги и постоянном
$I_{\mathrm{sp}}=300$~с. Предел $600$~с задан как учебное
ограничение, а не характеристика реального аппарата.}

\subsection{Отклонение уровня тяги}
\label{sec:thrust}

При постоянных тяге и удельном импульсе массовый расход
постоянен. Для того же $\Delta m$ длительность равна
\begin{equation}
t_e-t_b=\frac{c\Delta m}{F}.
\label{eq:duration}
\end{equation}
В предыдущем примере при тяге $450$~Н манёвр длится
$\OvDurationNominal$~с, а при $300$~Н ---
$\OvDurationLowThrust$~с.

Снижение тяги при том же $I_{\mathrm{sp}}$ и заданной
характеристической скорости увеличивает длительность,
но в модели \eqref{eq:fuel} не меняет расход.
Расход и длительность реагируют на разные характеристики ДУ.

Более долгий манёвр может выйти за тепловой предел,
потребовать иной ориентации или стать слишком протяжённым
относительно нужного орбитального положения.
На рисунке~\ref{fig:overview-propulsion},~б вариант $300$~Н
превышает установленный предел $600$~с.
Значит, одного увеличения длительности для продолжения
плана недостаточно.
\FloatBarrier

\subsection{Отказ двигателя и переход к резервным возможностям}
\label{sec:failure}

После отказа меняется множество доступных управлений.
Другой двигатель может иметь меньшую тягу или удельный импульс,
другое направление установки и другую систему подачи компонентов.
Нужно заново проверить последовательность, а не только заменить
одно значение тяги в прежнем расписании.

Например, основной двухкомпонентный двигатель использует горючее
и окислитель, а резервные однокомпонентные двигатели ---
только горючее. После отказа основного двигателя запас окислителя
может остаться на борту, но быть недоступным для резервного режима.
Наличие рабочего тела в баках ещё не определяет располагаемое
воздействие: важна возможность его использовать выбранной ДУ.

\begin{explanation}{Реальный пример: Akatsuki}
В 2010~году Akatsuki не выполнил запланированный выход
на орбиту Венеры после неисправности основной ДУ.
Для повторной попытки подготовили программу с использованием
двигателей ориентации~\cite{jaxa-plan}.
В декабре 2015~года они работали около 20~минут;
последующее определение орбиты подтвердило захват аппарата
Венерой~\cite{jaxa-burn,jaxa-result}.
\end{explanation}

Схема на рисунке~\ref{fig:overview-akatsuki} связывает отказ,
изменение плана и повторную попытку. Успешное продолжение
потребовало другого доступного управления, нового времени встречи
с Венерой и другой первоначальной орбиты захвата.

\coursefigure[.94]{overview-akatsuki}{Восстановление миссии Akatsuki:
нештатное событие 2010~года, подготовка продолжения и выход
на орбиту в 2015~году. Историческая схема по сообщениям JAXA;
она не изображает рассчитанную траекторию межпланетного перелёта.}

\subsection{Повторяемость и систематическое смещение}
\label{sec:repeatability}

Характеристики двигателя могут устойчиво отличаться от модели,
а могут меняться между манёврами. Первая ситуация создаёт
систематическое смещение. Вторая проявляется как неполная
повторяемость исполнения.

Причинами служат, например, различия давления и температуры,
теплового режима аппарата и состояния системы подачи.
Даже после уточнения среднего уровня воздействия
отдельные результаты не становятся одинаковыми.

На рисунке~\ref{fig:overview-repeatability} задан учебный ряд
из восьми манёвров: постоянный недобор $1.5\%$
сочетается с небольшими назначенными отклонениями.
После удаления известной постоянной составляющей разброс
сохраняется. Это иллюстрация смысла калибровки;
фактические данные и точность оценки коэффициента
в реальной миссии получают по измерениям.

\coursefigure[.92]{overview-repeatability}{Систематическое отклонение
и изменения между манёврами. Зелёная серия получена удалением
известного постоянного смещения из тех же учебных данных.
Она показывает сохраняющийся разброс, а не новые измерения
после реальной калибровки.}
\FloatBarrier

\subsection{Неопределённость соотношения компонентов}
\label{sec:mixture}

У двухкомпонентной ДУ имеет значение не только суммарная
заправленная масса. Обозначим через $K$ отношение расходуемой
массы окислителя к расходуемой массе горючего.
Если фактическое $K$ отличается от расчётного,
один компонент может закончиться раньше другого.

Пусть на борту $50$~кг горючего и $80$~кг окислителя.
При расчётном $K=1.6$ эти запасы согласованы.
При фактическом $K=1.8$ после расходования всего окислителя
будет использовано $80/1.8$~кг горючего.
Остаток составит
\begin{equation}
m_{\text{гор,ост}}=50-\frac{80}{1.8}
=\OvMixtureResidual~\text{кг}.
\label{eq:mixture}
\end{equation}

На рисунке~\ref{fig:overview-mixture} сопоставлены схема подачи
и такой остаток. Продолжить работу основного двухкомпонентного
двигателя за счёт одного горючего нельзя.
В показанной конфигурации оно доступно резервному
однокомпонентному двигателю, но его возможности будут другими.

\coursefigure[.94]{overview-mixture}{Доступность рабочего тела.
Схема показывает, какие компоненты использует каждый двигатель.
Сопоставление запасов иллюстрирует остаток горючего после
раннего исчерпания окислителя при $K=1.8$.
Все массы и схема ДУ заданы для учебного примера.}

Такой анализ нужен при заправке, оценке оставшегося ресурса
и выборе продолжения после отказа. «Остаток в баках»
и «рабочее тело, доступное для данного режима»
могут различаться.
\FloatBarrier

\section{Опоздание, пропуск и аварийное прекращение манёвра}
\label{sec:timing}

Эти ситуации меняют программу по-разному.
При опоздании воздействие всё же получено, но в другой момент.
При пропуске оно не получено в выбранном окне.
При аварийном прекращении воздействие выполнено частично.

\subsection{Опоздание: другая геометрия приложения импульса}
\label{sec:delay}

За время задержки аппарат продолжает движение.
Меняются положение, скорость и местные направления.
Поэтому прежние параметры манёвра могут дать другой результат;
позднее исполнение требует повторной оценки.

На рисунке~\ref{fig:overview-delay} рассмотрена учебная
эллиптическая орбита с $r_p=7000$~км и $r_a=12000$~км.
Номинальный трансверсальный импульс $100$~м/с в апогее
повышает перигей до $\TimingTargetRp$~км.
В этом сравнении целевыми условиями служат радиусы перигея
и апогея; плоскость сохраняется, направление линии апсид,
фаза и конечный срок свободны.

Если те же $100$~м/с приложены позже в текущем трансверсальном
направлении, изменяются оба радиуса апсид.
Для оценки последствий выберем конкретное продолжение:
в следующем перигее понизим апогей до $12000$~км,
а затем в апогее доведём перигей до целевого.
Суммарная характеристическая скорость $J$ включает
запоздалый импульс и оба исправляющих.

\coursefigure[.94]{overview-delay}{Цена позднего исполнения
в выбранном учебном сценарии. После задержанного импульса
использована указанная в тексте двухапсидальная коррекция
размеров орбиты. Показана добавка характеристической скорости
к номинальным $100$~м/с; это результат данной схемы продолжения.
Направление линии апсид и срок не ограничены.}

При задержке $120$~с добавка к $J$ составляет
$\TimingDelayShortCost$~м/с,
при $600$~с --- $\TimingDelayMediumCost$~м/с,
при $900$~с --- $\TimingDelayLongCost$~м/с.
График позволяет связать запас на коррекцию с допустимой
задержкой для выбранной схемы.

В подготовке операции возможны два подхода:
задать максимальную задержку и оценить необходимый резерв
либо задать резерв и найти допустимые задержки.
Перед выдачей поздней команды результат сравнивают
с другими вариантами, включая ожидание нового окна.

\subsection{Пропуск: воздействие не получено}
\label{sec:missed}

При полном пропуске аппарат продолжает движение
по предманёвренной орбите. Это меняет последующие события,
а иногда нарушает фазирование и условия применения следующих
манёвров. На рисунке~\ref{fig:overview-missed}
сопоставлены номинальное исполнение, опоздание и пропуск.

\coursefigure[.94]{overview-missed}{Временные последствия трёх
вариантов: номинальное исполнение, поздний импульс
с выбранным исправлением и ожидание следующего апогея.
Время относится к учебной модели со свободным сроком.}

В центральном поле для предыдущего примера можно дождаться
следующего апогея и выполнить прежние $100$~м/с.
Период исходной орбиты равен $\TimingInitialPeriod$~с.
Те же размеры целевой орбиты достигаются позже,
но суммарная характеристическая скорость остаётся $100$~м/с.

При задержке $120$~с выбранная выше схема исправления
завершается примерно через $9539$~с от номинального момента.
Ожидание следующего апогея заканчивается уже через
$9215$~с и требует меньшего $J$.
Следовательно, в данном сравнении перенос на новое окно
выгоднее исполнения запоздалого импульса с таким продолжением.

Это решение опирается на свободный срок и принятую модель.
При ограниченном времени, требованиях к фазе,
орбитальных возмущениях или изменившейся доступности
аппарата результат сравнения может быть другим.
Пропуск окна не означает универсального переноса всей программы
ровно на один виток.

\subsection{Аварийное прекращение: часть манёвра выполнена}
\label{sec:abort}

Причиной прекращения может быть выход температуры
или скорости вращения за допустимые пределы,
нарушение ограничений ориентации, неисправность двигателя
или ошибка программы. Аппарат оказывается на орбите,
которая зависит от момента прекращения.

На рисунке~\ref{fig:overview-abort} показан отдельный расчёт
протяжённого манёвра. Исходные радиусы апсид те же:
$7000$ и $12000$~км. Манёвр начинается в апогее;
тяга $450$~Н направлена по текущей трансверсали,
начальная масса $2000$~кг,
$I_{\mathrm{sp}}=300$~с.
Плановая длительность $\TimingBurnDuration$~с
соответствует характеристической скорости $100$~м/с.

\coursefigure[.94]{overview-abort}{Достигнутые орбиты при прекращении
протяжённого манёвра в разные моменты. Сравниваются оскулирующие
орбиты после прекращения тяги, а не синхронные положения аппаратов.
Опорный результат здесь --- полностью выполненный протяжённый
манёвр; он отличается от идеального импульсного результата.}

После половины планового времени получено
$\TimingAbortHalfDV$~м/с.
Доля времени не равна точно доле характеристической скорости:
масса уменьшается, а ускорение $F/m$ растёт.
Полностью выполненный протяжённый манёвр даёт
$r_p\approx7461.6$~км и $r_a\approx12002.9$~км.
Это собственный опорный результат данного расчёта.

После прекращения сначала восстанавливают безопасное состояние
КА и проверяют телеметрию. Для поиска и сопровождения аппарата
нужен первоначальный прогноз положения.
Затем уточняют достигнутую орбиту и полученное воздействие,
проверяют состояние ДУ и рассчитывают продолжение.

Заранее можно исследовать несколько моментов остановки,
как на рисунке, и подготовить варианты действий.
Полученное решение нельзя свести к исполнению остатка
старой длительности: аппарат уже движется по другой орбите,
а условия повторного запуска могут измениться.
\FloatBarrier

\section{Цикл орбитальных операций}
\label{sec:cycle}

Анализ отдельных ошибок становится полезным, когда встроен
в повторяющийся процесс управления. После каждого события
нужно получить данные, обновить модель, выбрать допустимое
продолжение и подготовить его исполнение.

\subsection{Оценка, прогноз, расчёт и исполнение}

На рисунке~\ref{fig:overview-cycle} стадии объединены
в пять крупных частей:
\begin{enumerate}
\item \textbf{Оценка состояния.}
Проверка аппарата, определение орбиты и ориентации,
реконструкция манёвра, уточнение модели ДУ.
\item \textbf{Прогноз событий.}
Сеансы связи, тень, ограничения датчиков и ориентации,
возможные сближения с другими объектами.
\item \textbf{Расчёт продолжения.}
Последовательность манёвров, анализ ошибок, проверка
ресурса и резервных сценариев.
\item \textbf{Подготовка и проверка команд.}
Преобразование параметров манёвра в команды конкретного КА,
проверка последовательности в имитационной модели.
\item \textbf{Согласование и контроль исполнения.}
Утверждение плана, загрузка команд, выполнение,
контроль телеметрии и переход к новой оценке состояния.
\end{enumerate}

\coursefigure[.80]{overview-cycle}{Повторяющийся цикл орбитальных
операций. Крупные стадии связывают оценку фактического состояния,
прогноз, подготовку управления и контроль результата.
Внутри показаны основные подэтапы~\cite[§15.4]{diprinzio}.}

Орбитальное состояние обновляют для сопровождения аппарата
и подготовки дальнейших сеансов. Прогноз событий строят
как для номинального исполнения, так и для предусмотренного
резервного варианта. Загрузка команд может требовать сеанса связи,
а само исполнение подготовленной программы бывает автономным.

\subsection{Реконструкция и калибровка в общем процессе}

Нужно различать команду, телеметрию работы двигателя
и изменение орбиты. Телеметрия подтверждает события
и помогает понять причины отклонения.
Орбитальные измерения позволяют оценить достигнутое движение.

Реконструкция оценивает выполненное воздействие,
сравнение с планом выявляет расхождение,
калибровка уточняет модель будущего исполнения~\cite[§15.3]{diprinzio}.
Скорости для импульсной реконструкции сравнивают
в одну эпоху и в одной системе координат.
Для протяжённого манёвра требуется отделить тяговое изменение
от движения под действием гравитации и возмущающих сил.

Результат имеет неопределённость. По одному восстановленному
модулю импульса нельзя независимо определить тягу,
массу, расход и удельный импульс.
Полезно совместно рассматривать орбитальные данные,
телеметрию и последовательность манёвров.

\subsection{Операционное решение при пониженной тяге}

Вернёмся к примеру из раздела~\ref{sec:thrust}.
Перед предстоящим манёвром текущая масса КА равна $2000$~кг.
Пусть по имеющимся данным модель уточнена:
доступная тяга равна $300$ вместо $450$~Н,
а $I_{\mathrm{sp}}$ остался $300$~с.
Для следующих $100$~м/с потребуется
$\OvDurationLowThrust$~с, но допустимая длительность ---
только $600$~с.

Если ограничиться $600$~с работы, аппарат получит
$\OvLimitedDV$~м/с. Требование не будет выполнено.
Теперь нужно проверить другие варианты:
разделение воздействия между несколькими манёврами
с допустимыми паузами, другую конфигурацию двигателей
или изменение программы перехода.

Каждый вариант проходит полный цикл:
расчёт орбиты и событий, проверку ограничений,
анализ ошибок и ресурса, подготовку команд и контроль.
Так сведения об ошибке ДУ превращаются в конкретное
изменение операции.

\begin{explanation}{Условие принятия нового плана}
Продолжение должно достигать приемлемой цели
доступными средствами, укладываться в ограничения и ресурс.
После исполнения его результат снова проверяется по данным.
\end{explanation}
\FloatBarrier

\begin{thebibliography}{9}
\addcontentsline{toc}{section}{Литература}
\bibitem{diprinzio}
DiPrinzio M. \textit{Methods of Orbital Maneuvering}.
American Institute of Aeronautics and Astronautics, 2019.
Гл.~15: ошибки выведения, наведения и ДУ (§15.1.1, 15.1.3--15.1.4),
опоздание, пропуск и прерывание (§15.1.5--15.1.7),
реконструкция и калибровка (§15.3), цикл операций (§15.4).

\bibitem{esa-recovery}
ESA. \textit{Galileo satellite recovered and transmitting navigation signals}.
3 December 2014.
\href{https://www.esa.int/Newsroom/Press_Releases/Galileo_satellite_recovered_and_transmitting_navigation_signals}
{Сообщение о восстановлении пятого спутника Galileo}.

\bibitem{esa-facts}
ESA. \textit{Galileo: Facts and figures}.
\href{https://www.esa.int/Applications/Satellite_navigation/Galileo/Facts_and_figures}
{Справочные размеры штатной и исправленной орбит}.
Дата обращения: 17~ноября 2026~года.

\bibitem{jaxa-plan}
JAXA. \textit{Akatsuki: повторная попытка выхода на орбиту Венеры}.
6 February 2015.
\href{https://www.isas.jaxa.jp/topics/files/20150206_akatsuki.pdf}
{Материалы JAXA о плане применения двигателей ориентации
вместо неисправного основного двигателя}.

\bibitem{jaxa-burn}
JAXA. \textit{Результат работы двигателей ориентации
при выходе Akatsuki на орбиту Венеры}. 7 December 2015.
\href{https://www.jaxa.jp/press/2015/12/20151207_akatsuki_j.html}
{Сообщение JAXA о выполнении манёвра}.

\bibitem{jaxa-result}
JAXA. \textit{Venus Climate Orbiter «AKATSUKI» Inserted Into Venus' Orbit}.
9 December 2015.
\href{https://global.jaxa.jp/press/2015/12/20151209_akatsuki.html}
{Подтверждение выхода на орбиту по результатам её определения}.
\end{thebibliography}
\end{document}
